\documentclass{exam-zh}
\usepackage{edu-explanation}

\examsetup{
  question/show-answer = true,
  fillin/show-answer = true,
  solution/show-solution = show-stay,
}

\begin{document}

\eduSectionTitle{例题 1：已知全长，求长段}

\begin{eduProblemBlock}[例题 1]
一条装饰线全长 $30\text{ cm}$，按黄金分割分成较长、较短两段。
求较长一段的长度。

\end{eduProblemBlock}





\begin{eduAuxItem}{思路导航}
列式计算\end{eduAuxItem}




\begin{eduSubQuestionItem}{例题 1}
求较长一段的长度。\end{eduSubQuestionItem}

\begin{eduDualExplain}
  \begin{eduExplainLeft}
    \eduColumnTitle{\eduIconBulb}{小贴士}
        \eduSideHint{第一步：确定两条边}{本题已知的是“全” $30\text{ cm}$，要求的是“长”。
}
        \eduSideHint{第二步：确定比值}{三组两两比值为
\[
  \begin{aligned}
    \frac{\text{短}}{\text{长}}&=\frac{\sqrt5-1}{2},\\
    \frac{\text{长}}{\text{全}}&=\frac{\sqrt5-1}{2},\\
    \frac{\text{短}}{\text{全}}&=\frac{3-\sqrt5}{2}.
  \end{aligned}
\]
本题用“长比全” $\dfrac{\text{长}}{\text{全}}=\dfrac{\sqrt5-1}{2}$。
}
        \eduSideHint{第三步：确定乘除}{已知的是分母“全”，所以用 $30\times\dfrac{\sqrt5-1}{2}$。
}
  \end{eduExplainLeft}

  \begin{eduExplainRight}
    \eduColumnTitle{\eduIconBook}{解答}
      \eduExplainStep{1}{列式计算}{\[
  30\times\frac{\sqrt5-1}{2}
  =15(\sqrt5-1)\text{ cm}.
\]
}
  \end{eduExplainRight}
\end{eduDualExplain}



\clearpage
\eduSectionTitle{例题 2：已知短段，求全长}

\begin{eduProblemBlock}[例题 2]
一段窗格按黄金分割分成长、短两段，其中较短一段长 $6\text{ m}$。
求整段窗格的长度。

\end{eduProblemBlock}





\begin{eduAuxItem}{思路导航}
列式计算\end{eduAuxItem}




\begin{eduSubQuestionItem}{例题 2}
求整段窗格的长度。\end{eduSubQuestionItem}

\begin{eduDualExplain}
  \begin{eduExplainLeft}
    \eduColumnTitle{\eduIconBulb}{小贴士}
        \eduSideHint{第一步：确定两条边}{本题已知的是“短” $6\text{ m}$，要求的是“全”。
}
        \eduSideHint{第二步：确定比值}{直接使用“短比全” $\dfrac{\text{短}}{\text{全}}=\dfrac{3-\sqrt5}{2}$，不必先求长段。
}
        \eduSideHint{第三步：确定乘除}{已知的是分子“短”，所以用 $6\div\dfrac{3-\sqrt5}{2}$。
}
  \end{eduExplainLeft}

  \begin{eduExplainRight}
    \eduColumnTitle{\eduIconBook}{解答}
      \eduExplainStep{1}{列式计算}{\[
  6\div\frac{3-\sqrt5}{2}
  =3(3+\sqrt5)\text{ m}.
\]
}
  \end{eduExplainRight}
\end{eduDualExplain}




\end{document}
